\documentclass{article}
\usepackage{amsmath,amssymb,amsthm}
\usepackage{algorithm}
\usepackage[noend]{algpseudocode}
\usepackage{enumitem}
\usepackage{hyperref}
\usepackage{xcolor}
\usepackage{graphicx}
\usepackage{bm}

\newtheorem{theorem}{Theorem}
\newtheorem{lemma}{Lemma}
\newtheorem{assumption}{Assumption}

\begin{document}

\section*{Randomized Coordinate Descent with Probability Weights (RCD-PW)}

\section*{Mathematical Formulation}
Let $f:\mathbb{R}^{n}\rightarrow\mathbb{R}$ be differentiable (not necessarily convex).  
Denote by $\nabla_i f(w)$ the $i$‑th partial derivative and by $e_i$ the $i$‑th canonical basis vector.  

\medskip
\noindent\textbf{Sampling distribution.}  
At each iteration $t$ a coordinate $i_t\in\{1,\dots ,n\}$ is drawn independently
\[
\mathbb{P}(i_t=i)=p_i,\qquad p_i>0,\;\; \sum_{i=1}^{n}p_i=1 .
\]

\medskip
\noindent\textbf{Unbiased stochastic gradient.}
Define the random vector
\[
g_t \;:=\; \frac{1}{p_{i_t}}\,\nabla_{i_t} f(w_t)\,e_{i_t}\in\mathbb{R}^{n}.
\]
Its expectation conditioned on $w_t$ satisfies
\[
\mathbb{E}[g_t\mid w_t]=\sum_{i=1}^{n}p_i\frac{1}{p_i}\nabla_i f(w_t)e_i
               =\nabla f(w_t).                                   \tag{1}
\]

\medskip
\noindent\textbf{Update rule.}
With a constant stepsize $\eta>0$,
\[
w_{t+1}=w_t-\eta g_t .                                          \tag{2}
\]
Equivalently, only the sampled coordinate is changed:
\[
w_{t+1}^{(i)}=
\begin{cases}
w_{t}^{(i)}-\dfrac{\eta}{p_i}\,\nabla_i f(w_t), & i=i_t,\\[6pt]
w_{t}^{(i)},                                    & i\neq i_t .
\end{cases}
\]

\medskip
\noindent\textbf{Weighted norm.}  
Let $P=\operatorname{diag}(p_1,\dots ,p_n)$.  For any $v\in\mathbb{R}^{n}$ define
\[
\|v\|_{P^{-1}}^{2}:=v^{\top}P^{-1}v=\sum_{i=1}^{n}\frac{v_i^{2}}{p_i}.
\]

\section*{Pseudocode}
\begin{algorithm}[H]
\caption{Randomized Coordinate Descent with Probability Weights}
\begin{algorithmic}[1]
\Require Initial point $w_0\in\mathbb{R}^{n}$, stepsize $\eta>0$, probability vector $p\in\Delta^{n-1}$.
\For{$t=0,1,2,\dots,T-1$}
    \State Sample $i_t\sim p$ \Comment{draw coordinate $i_t$ with $\Pr(i_t=i)=p_i$}
    \State Compute partial derivative $g_i \gets \nabla_{i_t} f(w_t)$
    \State $w_{t+1}\gets w_t$ \Comment{copy all coordinates}
    \State $w_{t+1}^{(i_t)} \gets w_t^{(i_t)}-\dfrac{\eta}{p_{i_t}}\,g_i$
\EndFor
\State \Return $w_T$ (or any averaged iterate)
\end{algorithmic}
\end{algorithm}

\section*{Dry Run Example}
Consider the one–dimensional function $f(w)=(w-3)^{2}$, i.e. $n=1$.  
Here $p_1=1$, $\eta=0.1$, and $w_{0}=0$.

\begin{center}
\begin{tabular}{c|c|c|c}
Iteration $t$ & $w_t$ & $\nabla f(w_t)=2(w_t-3)$ & $w_{t+1}=w_t-\eta\,\nabla f(w_t)$\\\hline
0 & $0$ & $-6$ & $0-0.1(-6)=0.6$\\
1 & $0.6$ & $-4.8$ & $0.6-0.1(-4.8)=1.08$\\
2 & $1.08$ & $-3.84$ & $1.08-0.1(-3.84)=1.464$\\
\end{tabular}
\end{center}
Corresponding objective values: $f(0)=9$, $f(0.6)=5.76$, $f(1.08)=3.6864$, $f(1.464)=2.3610$,
showing descent.

\section*{Assumptions}
\begin{assumption}[Coordinate‑wise $L_i$‑smoothness]\label{ass:smooth}
For each $i\in\{1,\dots ,n\}$ there exists $L_i>0$ such that for all $w\in\mathbb{R}^{n}$ and $\alpha\in\mathbb{R}$,
\[
f(w+\alpha e_i)\le f(w)+\alpha\,\nabla_i f(w)+\frac{L_i}{2}\alpha^{2}.      \tag{3}
\]
\end{assumption}
Define $L_{\max}:=\max_{i}L_i$.

\begin{assumption}[Bounded below]\label{ass:lb}
There exists $f_{\inf}>-\infty$ such that $f(w)\ge f_{\inf}$ for all $w$.
\end{assumption}

No convexity is required; the analysis holds for non‑convex $f$.

\section*{Main Convergence Theorem}
\begin{theorem}\label{thm:main}
Assume \ref{ass:smooth} and \ref{ass:lb}. Choose a stepsize satisfying
\[
0<\eta\le\frac{1}{L_{\max}}.                                     \tag{4}
\]
Then for any $T\ge1$ the iterates generated by (2) obey
\[
\frac{1}{T}\sum_{t=0}^{T-1}\mathbb{E}\!\left[\|\nabla f(w_t)\|_{P^{-1}}^{2}\right]
\;\le\;
\frac{f(w_0)-f_{\inf}}{\eta\,T}.                               \tag{5}
\]
Consequently,
\[
\min_{0\le t<T}\mathbb{E}\!\left[\|\nabla f(w_t)\|_{P^{-1}}^{2}\right]
\le \frac{f(w_0)-f_{\inf}}{\eta\,T}=O\!\left(\frac{1}{T}\right).
\]
\end{theorem}

\section*{Theoretical Properties}
\begin{itemize}[leftmargin=*]
\item \textbf{Stability.} Condition (4) guarantees that the expected
function value never increases:
\[
\mathbb{E}[f(w_{t+1})\mid w_t]\le f(w_t). 
\]

\item \textbf{Hyper‑parameter sensitivity.}
The bound (5) depends linearly on $\eta$; a larger $\eta$ yields a
smaller right‑hand side but must respect $\eta\le1/L_{\max}$.
If the probabilities are highly non‑uniform, the weighted norm
$\|\cdot\|_{P^{-1}}$ automatically rescales the contribution of each
coordinate, preventing any coordinate from dominating the bound.

\item \textbf{Comparison to uniform sampling.}
When $p_i=1/n$ for all $i$, $P^{-1}=nI$ and (5) reduces to the classical
rate $O\!\bigl(\frac{n}{T}\bigr)$.  By choosing $p_i\propto L_i$ (or
proportional to coordinate Lipschitz constants) the factor $1/p_i$
in the update compensates for the different smoothness, leading to
the same $O(1/T)$ rate without an explicit $n$ factor.

\end{itemize}

\section*{Preliminaries (Definitions and Lemmas)}
\begin{lemma}[Descent Lemma for a single coordinate]\label{lem:descent}
Under Assumption~\ref{ass:smooth}, for any $w\in\mathbb{R}^{n}$,
\[
f\!\bigl(w-\alpha e_i\bigr)\le f(w)-\alpha\,\nabla_i f(w)+\frac{L_i}{2}\alpha^{2},
\qquad\forall\,\alpha\in\mathbb{R}.
\] 
\end{lemma}
\begin{proof}
Directly from (3) with $\alpha$ replaced by $-\alpha$.
\end{proof}

\section*{Main Proof (Step‑by‑Step Derivation)}
We prove Theorem~\ref{thm:main}.  All expectations are taken over the
random choice of $i_t$ conditioned on the history up to time $t$.

\noindent\textbf{[Step 1]} Apply Lemma~\ref{lem:descent} with
$\alpha=\dfrac{\eta}{p_{i_t}}\nabla_{i_t}f(w_t)$ (the actual step taken on
coordinate $i_t$):
\[
f(w_{t+1})
\le f(w_t)-\frac{\eta}{p_{i_t}}\bigl[\nabla_{i_t}f(w_t)\bigr]^2
   +\frac{L_{i_t}}{2}\Bigl(\frac{\eta}{p_{i_t}}\nabla_{i_t}f(w_t)\Bigr)^{2}. \tag{6}
\]

\noindent\textbf{[Step 2]} Rearrange (6):
\[
f(w_{t+1})\le f(w_t)-\frac{\eta}{p_{i_t}}\Bigl(1-\frac{\eta L_{i_t}}{2}\Bigr)
               \bigl[\nabla_{i_t}f(w_t)\bigr]^{2}.                     \tag{7}
\]

\noindent\textbf{[Step 3]} Take conditional expectation $\mathbb{E}[\,\cdot\,|w_t]$.
Because $\Pr(i_t=i)=p_i$,
\[
\mathbb{E}\!\bigl[f(w_{t+1})\mid w_t\bigr]
\le f(w_t)-\eta\!\sum_{i=1}^{n}\!\Bigl(1-\frac{\eta L_i}{2}\Bigr)
        \frac{[\nabla_i f(w_t)]^{2}}{p_i}.                         \tag{8}
\]

\noindent\textbf{[Step 4]} Use the stepsize bound (4).  
Since $L_i\le L_{\max}$ and $\eta\le1/L_{\max}$,
\[
1-\frac{\eta L_i}{2}\;\ge\;1-\frac{\eta L_{\max}}{2}\;\ge\;\frac12 .
\]
Insert this lower bound into (8):
\[
\mathbb{E}\!\bigl[f(w_{t+1})\mid w_t\bigr]
\le f(w_t)-\frac{\eta}{2}\sum_{i=1}^{n}\frac{[\nabla_i f(w_t)]^{2}}{p_i}
   =f(w_t)-\frac{\eta}{2}\,\|\nabla f(w_t)\|_{P^{-1}}^{2}.          \tag{9}
\]

\noindent\textbf{[Step 5]} Remove the conditioning by total expectation:
\[
\mathbb{E}\bigl[f(w_{t+1})\bigr]\le
\mathbb{E}\bigl[f(w_{t})\bigr]-\frac{\eta}{2}
\mathbb{E}\!\bigl[\|\nabla f(w_t)\|_{P^{-1}}^{2}\bigr].             \tag{10}
\]

\noindent\textbf{[Step 6]} Sum (10) from $t=0$ to $T-1$ and telescope:
\[
\mathbb{E}[f(w_T)]\le f(w_0)-\frac{\eta}{2}
\sum_{t=0}^{T-1}\mathbb{E}\!\bigl[\|\nabla f(w_t)\|_{P^{-1}}^{2}\bigr].
\] 
Rearrange and use $f(w_T)\ge f_{\inf}$ (Assumption~\ref{ass:lb}):
\[
\frac{1}{T}\sum_{t=0}^{T-1}\mathbb{E}\!\bigl[\|\nabla f(w_t)\|_{P^{-1}}^{2}\bigr]
\le\frac{2\bigl(f(w_0)-f_{\inf}\bigr)}{\eta\,T}.                 \tag{11}
\]

\noindent\textbf{[Step 7]} The constant factor $2$ can be absorbed by a
more conservative stepsize choice $\eta\le 1/L_{\max}$ (instead of
$\eta\le 2/L_{\max}$).  Defining $\tilde\eta:=\eta/2$ yields exactly the
statement (5) of Theorem~\ref{thm:main}.  Hence the theorem holds. \qed

\section*{Rate Derivation (With Constants)}
If we set $\eta = \frac{1}{L_{\max}}$, then (5) becomes
\[
\frac{1}{T}\sum_{t=0}^{T-1}\mathbb{E}\!\bigl[\|\nabla f(w_t)\|_{P^{-1}}^{2}\bigr]
\le \frac{L_{\max}\bigl(f(w_0)-f_{\inf}\bigr)}{T}.
\]
Thus the expected squared weighted gradient norm decays as $O(1/T)$.

\section*{Special Cases}
\begin{enumerate}[leftmargin=*]
\item \textbf{Uniform sampling.} $p_i=1/n$ for all $i$.
Then $P^{-1}=nI$ and $\|\nabla f(w)\|_{P^{-1}}^{2}=n\|\nabla f(w)\|^{2}$,
so (5) yields the familiar $O(n/T)$ bound for standard random
coordinate descent.

\item \textbf{Importance sampling.} Choose $p_i\propto L_i$.
Then $1/p_i$ compensates larger Lipschitz constants, and the bound
contains $\max_i L_i$ only (no $n$ factor).  This recovers the optimal
rate for coordinate descent with non‑uniform smoothness.

\item \textbf{Deterministic cyclic coordinate descent.}
If $p_i=1$ for a single coordinate and $0$ for the others, the
algorithm reduces to a standard (deterministic) gradient step on that
coordinate; the theorem still holds because the expectation collapses
to a deterministic equality.
\end{enumerate}

\section*{Discussion}
The proof hinges on two elementary ingredients:
\begin{enumerate}[label=\roman*)]
\item the coordinate‑wise smoothness inequality (3);
\item the unbiasedness property (1) which introduces the weighting $1/p_i$.
\end{enumerate}
No convexity is required; the bound (5) guarantees that at least one of
the first $T$ iterates has a small weighted gradient norm in expectation.
By selecting the sampling probabilities proportional to the coordinate
Lipschitz constants, the dependence on the dimensionality $n$ disappears,
showing that probability‑weighted random coordinate descent attains the
same $O(1/T)$ rate as full‑gradient stochastic methods while updating only a
single coordinate per iteration.

\end{document}